\documentclass[12pt]{article}
\usepackage[utf8]{inputenc}
\usepackage{amsmath, amssymb, amsthm}
\usepackage{booktabs}
\usepackage{geometry}
\geometry{a4paper, margin=1in}

\newtheorem{theorem}{Theorem}
\newtheorem{lemma}[theorem]{Lemma}
\newtheorem{corollary}[theorem]{Corollary}
\newtheorem{proposition}[theorem]{Proposition}
\theoremstyle{definition}
\newtheorem{definition}[theorem]{Definition}

\title{Problem 54: Poker Hands}
\author{Project Euler Solution}
\date{}

\begin{document}
\maketitle

\section{Problem Statement}
Given 1000 poker hands dealt to two players (from \texttt{poker.txt}), determine how many hands Player~1 wins.

\section{Formal Development}

\begin{definition}[Card and Hand]
A \emph{card} is a pair $(v, s)$ with value $v \in \{2, 3, \ldots, 14\}$ and suit $s \in \{\clubsuit, \diamondsuit, \heartsuit, \spadesuit\}$. A \emph{hand} is an ordered 5-tuple of distinct cards.
\end{definition}

\begin{definition}[Derived Quantities]
For a hand $H$ with values $v_1, \ldots, v_5$ and suits $s_1, \ldots, s_5$:
\begin{itemize}
    \item \emph{Frequency}: $f(v) = |\{i : v_i = v\}|$.
    \item \emph{Flush}: $\phi(H) \iff s_1 = \cdots = s_5$.
    \item \emph{Straight}: $\sigma(H) \iff$ the sorted distinct values are 5 consecutive integers, or $\{2,3,4,5,14\}$.
\end{itemize}
\end{definition}

\begin{definition}[Hand Category]
The category $c(H) \in \{0, \ldots, 9\}$:
\begin{center}
\begin{tabular}{cll}
\toprule
$c$ & Name & Condition \\
\midrule
9 & Royal Flush & $\sigma \wedge \phi \wedge \{v_i\} = \{10,11,12,13,14\}$ \\
8 & Straight Flush & $\sigma \wedge \phi \wedge \neg(c=9)$ \\
7 & Four of a Kind & Freq.\ partition $= (4,1)$ \\
6 & Full House & Freq.\ partition $= (3,2)$ \\
5 & Flush & $\phi \wedge \neg\sigma$ \\
4 & Straight & $\sigma \wedge \neg\phi$ \\
3 & Three of a Kind & Freq.\ partition $= (3,1,1)$ \\
2 & Two Pairs & Freq.\ partition $= (2,2,1)$ \\
1 & One Pair & Freq.\ partition $= (2,1,1,1)$ \\
0 & High Card & otherwise \\
\bottomrule
\end{tabular}
\end{center}
\end{definition}

\begin{definition}[Comparison Key]
The comparison key is:
\[
\kappa(H) = \bigl(c(H),\; \tau_1, \tau_2, \ldots, \tau_m\bigr)
\]
where $(\tau_i)$ are the tiebreaker values: for straights/straight flushes, the high card of the straight (5 for the wheel); for all other categories, distinct values sorted by $(f(v), v)$ descending.
\end{definition}

\begin{theorem}[Correctness]
\label{thm:correct}
For two hands $H_1, H_2$ with no ties: $H_1$ beats $H_2 \iff \kappa(H_1) > \kappa(H_2)$ lexicographically.
\end{theorem}

\begin{proof}
The category $c$ is the primary discriminator. Within each category, the tiebreaker tuple orders groups by frequency (then value) descending, matching standard poker rules:
\begin{itemize}
    \item \textbf{Four of a Kind}: compare quad value, then kicker.
    \item \textbf{Full House}: compare triple value, then pair value.
    \item \textbf{Two Pairs}: compare higher pair, lower pair, then kicker.
    \item \textbf{One Pair}: compare pair value, then kickers descending.
    \item \textbf{High Card / Flush}: compare all values descending.
    \item \textbf{Straight types}: compare by high card.
\end{itemize}
In every case the lexicographic order on $\kappa$ faithfully implements the poker comparison.
\end{proof}

\begin{lemma}[Ace-Low Straight]
When sorted values are $\{2,3,4,5,14\}$, the ace functions as~1. The high card is~5.
\end{lemma}

\begin{proof}
By poker rules, $A$-$2$-$3$-$4$-$5$ is the lowest straight. Without this correction, the ace (14) would rank this straight above $2$-$3$-$4$-$5$-$6$.
\end{proof}

\section{Algorithm}

\begin{enumerate}
    \item For each of the 1000 lines:
    \begin{enumerate}
        \item Parse 10 cards; assign first 5 to $H_1$, last 5 to $H_2$.
        \item Compute $\kappa(H_1)$ and $\kappa(H_2)$.
        \item If $\kappa(H_1) > \kappa(H_2)$: increment win counter.
    \end{enumerate}
    \item Return win counter.
\end{enumerate}

\section{Complexity}

\begin{proposition}[Time]
$O(H)$ where $H = 1000$. Each hand evaluation is $O(1)$ (5 cards, fixed-size operations).
\end{proposition}

\begin{proposition}[Space]
$O(1)$ per hand evaluation.
\end{proposition}

\section{Result}
Player~1 wins $\boxed{376}$ hands out of 1000.

\section{Answer}

\[
\boxed{376}
\]

\end{document}
